\chapter*{Véc-tơ}

\section{Bài tập tại lớp}

\ 

\exercise Cho một điện tích $q > 0$ chuyển động trong vùng có từ trường $\vec{B}(0, 0, B)$ với vận tốc $\vec{v}(v_x, v_y, v_z)$. Xác định lực Lo-ren-xơ tác dụng lên điện tích $q$. Biết công thức lực Lo-ren-xơ là 
$$
\vec{F}_L = q(\vec{v} \wedge \vec{B}).
$$

\solution

\begin{align*}
   \vec{F}_L &= q(\vec{v} \wedge \vec{B}) \\
   &= q\left(\left(v_x \vec{i} + v_y \vec{j} + v_z \vec{k}\right) \wedge \left(B\vec{k}\right)\right) \\
   &= q\left(v_x B \left(\vec{i} \wedge \vec{k}\right) + v_y B \left(\vec{j} \wedge \vec{k}\right) + v_z B \left(\vec{k} \wedge \vec{k}\right)\right) \\
   &= q\left(v_x B \left(-\vec{j}\right) + v_y B \vec{i} + v_z B \vec{0}\right) \\
   &= q v_y B \vec{i} - q v_x B \vec{j}.
\end{align*}

\exercise Một vật có khối lượng $m = \qty{2}{\kg}$ được đặt đứng yên trên một mặt phẳng nghiêng nhám có hệ số ma sát $\mu = \num{0,4}$. 

\begin{enumerate}
   \item Ban đầu góc nghiêng $\alpha$ nhỏ nên vật vẫn cân bằng. Hỏi tăng chậm dần góc $\alpha$ đến giá trị $\alpha_0$ bằng bao nhiêu thì vật bắt đầu trượt?
   \item Cho góc nghiêng $\alpha = 30^\circ$. Vật trượt trên mặt phẳng nghiêng một đoạn bằng \qty{5}{m}. Xác định công của từng ngoại lực tác dụng lên vật.
\end{enumerate}

\solution

\begin{figure}[H]
   \centering
   \begin{tikzpicture}[>=latex]
      % Đổi tên biến góc thành \goc để tránh đè lên lệnh \alpha của hệ thống
      \def\goc{22}
      
      % Vẽ mặt đất (nằm ngang)
      \draw[thick] (-0.5, 0) -- (6, 0);
      
      % Vẽ mặt phẳng nghiêng
      \draw[thick] (0, 0) -- (\goc:5.8);
      
      % Vẽ ký hiệu góc alpha (hiện chữ alpha thay vì số 22)
      \draw[thick] (1, 0) arc (0:\goc:1);
      \node at ({\goc/2}:1.3) {$\alpha$};
      
      % Vẽ vật và các lực tác dụng
      \begin{scope}[shift={(\goc:3)}, rotate=\goc]
         % Khối hộp
         \draw[very thick, fill=white] (-0.6, 0) rectangle (0.6, 0.8);
         
         % Tâm vật
         \coordinate (CoM) at (0, 0.4);
         
         \draw[dashed, colorDodgerblue, ->] (-2.2, 0.4) -- (1.6, 0.4) node[right] {$x$};
         % Trục Oy hướng vuông góc lên trên mặt phẳng nghiêng
         \draw[dashed, colorDodgerblue, ->] (0, -1.0) -- (0, 2.2) node[above] {$y$};
         % Ký hiệu gốc tọa độ O
         \node[above right, colorDodgerblue, scale=0.8] at (CoM) {$O$};

         % Lực ma sát nghỉ
         \draw[colorCrimson, very thick, ->] (0, 0) -- (0.8, 0) node[below right] {$\vec{F}_{\text{ms}}$};
         
         % Phản lực N
         \draw[colorCrimson, very thick, ->] (CoM) -- (0, 1.6) node[above left] {$\vec{N}$};
         
         % Trọng lực P hướng thẳng đứng xuống dưới
         \draw[colorCrimson, very thick, ->] (CoM) -- ++({-90-\goc}:1.6) node[below] {$\vec{P}$};
      \end{scope}
   \end{tikzpicture}
\end{figure}

\setcounter{subexercise}{1}
\arabic{subexercise}. Khi vật đứng yên cân bằng trên mặt phẳng nghiêng,  hợp lực tác dụng lên vật bằng $\vec{0}$:
$$\vec{P} + \vec{N} + \vec{F}_{\text{ms}} = \vec{0}.$$

Chiếu lên hai trục tọa độ $Ox$ và $Oy$, ta lại có
\begin{equation*}
   \begin{cases}
      \vec{P} = - P \cdot \sin\alpha \vec{i} -  P \cdot \cos\alpha \vec{j} = - 10m\sin\alpha \vec{i} -  10m\cos\alpha \vec{j} \\
      \vec{N} = N \vec{j} \\
      \vec{F}_{\text{ms}} = F_{\text{ms}} \vec{i}
   \end{cases}.
\end{equation*}
Thay vào phương trình cân bằng để có:
\begin{align*}
   - 10m\sin\alpha \vec{i} -  10m\cos\alpha \vec{j} + N \vec{j} + F_{\text{ms}} \vec{i} &= \vec{0}; \\
   \left(F_{\text{ms}} - 10m\sin\alpha\right) \vec{i} + (N - 10m\cos\alpha) \vec{j} &= 0\vec{i} + 0\vec{j}.
\end{align*}
Cân bằng hệ số:
\begin{align*}
& \begin{cases}
N - (10 \cdot m)\cos\alpha = 0 \\
F_{\text{ms}} - 10 \cdot m\sin\alpha = 0
\end{cases} \\
& \begin{cases}
N = 10m \cdot \cos\alpha \\
F_{\text{ms}} = 10m \cdot \sin\alpha
\end{cases}.
\end{align*}
Mà ta lại có để vật không bị trượt thì:
\begin{equation*}
F_{\text{ms}} \le \mu \cdot N,
\end{equation*}
do vậy
\begin{align*}
10m \cdot \sin\alpha &\le \mu \cdot 10m \cdot \cos\alpha \\
\tan\alpha &\le \mu \\
\alpha &\le \arctan\mu \\
\alpha &\le \num{21.8}^\circ.
\end{align*}
Vậy góc tới hạn để vật không trượt là \boxed{\alpha_0 = \num{21.8}^\circ}.

\stepcounter{subexercise}
\arabic{subexercise}. Khi góc nghiêng $\alpha = 30^\circ > \alpha_0$, vật trượt xuống dưới dọc theo mặt phẳng nghiêng. Vì chiều dương của trục $Ox$ hướng lên trên nên
độ dịch chuyển $\vec{s}$ ngược chiều với vectơ đơn vị $\vec{i}$:
$$\vec{s} = -5 \vec{i}.$$

Biểu diễn các lực:
\begin{equation*} \begin{cases} \vec{P} = - 10m\sin\alpha \vec{i}
   - 10m\cos\alpha \vec{j} = -10 \vec{i} - 10\sqrt{3} \vec{j} \\ 
   \vec{N} = 10m\cos\alpha \vec{j} = 10\sqrt{3} \vec{j}\\ 
   \vec{F}_{\text{ms}}
   = \mu N \vec{i} = 10\mu m\cos\alpha \vec{i} = 4\sqrt{3} \vec{i}
   \end{cases}. 
\end{equation*}

Công của từng ngoại lực tác dụng lên vật: 
\begin{itemize} 
   \item Công của trọng lực $\vec{P}$: 
   \begin{equation*} 
      A_P = \vec{P} \cdot \vec{s} = \left(-10 \vec{i} - 10\sqrt{3} \vec{j}\right) \cdot \left(-5 \vec{i}\right) = (-10) \cdot (-5) = \boxed{\qty{50}{J}}. 
   \end{equation*}

   \item Công của phản lực $\vec{N}$: 
   \begin{equation*} 
      A_N = \vec{N} \cdot \vec{s} = \left(10\sqrt{3} \vec{j}\right) \cdot \left(-5 \vec{i}\right) = \boxed{\qty{0}{J}}. 
   \end{equation*}

   \item Công của lực ma sát trượt $\vec{F}_{\text{ms}}$: 
   \begin{equation*} 
      A_{\text{ms}} = \vec{F}_{\text{ms}} \cdot \vec{s} = \left(4\sqrt{3} \vec{i}\right) \cdot \left(-5 \vec{i}\right) = -20\sqrt{3} \text{ J} \approx \boxed{\qty{-34.6}{ J}}. 
   \end{equation*} 
\end{itemize}

\exercise Biết $F_1 = \qty{20}{N}$, $F_2 = \qty{30}{N}$, $F_3 = \qty{40}{N}$. Tìm độ lớn hợp lực $$\vec{F} = \vec{F}_1 + \vec{F}_2 + \vec{F}_3.$$

\solution
\begin{figure}[H]
   \centering
   \begin{tikzpicture}[>=latex]
      % Gốc tọa độ O của hệ quy chiếu phụ (góc trên bên trái)
      \draw[dashed, color=colorDodgerblue, ->] (-2.5, 1.2) -- (-1.8, 1.2) node[below] {$x$};
      \draw[dashed, color=colorDodgerblue, ->] (-2.5, 1.2) -- (-2.5, 1.9) node[left] {$y$};
      \node[color=colorDodgerblue, below left] at (-2.5, 1.2) {$O$};

      % Điểm đồng quy
      \coordinate (O) at (0, 0);
      \fill (O) circle (\pointSize);

      % Trục nằm ngang nét đứt
      \draw[dashed] (-2.2, 0) -- (2.2, 0);

      % Vẽ các vectơ lực (độ dài vẽ tỉ lệ thuận với độ lớn thực tế)
      % F1 = 20 N (vẽ dài 1 cm), F2 = 30 N (vẽ dài 1.5 cm), F3 = 40 N (vẽ dài 2.0 cm)
      \draw[very thick, ->] (O) -- (90:1.0) node[right] {$\vec{F}_1$};
      \draw[very thick, ->] (O) -- (330:1.5) node[right] {$\vec{F}_2$};
      \draw[very thick, ->] (O) -- (210:2.0) node[below left] {$\vec{F}_3$};

      % Vẽ các cung góc 120 độ
      % Góc giữa F1 và F2
      \draw (90:0.4) arc (90:-30:0.4);
      \node at (30:0.6) {$120^\circ$};

      % Góc giữa F2 và F3
      \draw (330:0.4) arc (330:210:0.4);
      \node at (270:0.6) {$120^\circ$};

      % Góc giữa F3 và F1
      \draw (210:0.4) arc (210:90:0.4);
      \node at (150:0.6) {$120^\circ$};
   \end{tikzpicture}
\end{figure}


Chọn hệ trục tọa độ vuông góc $Oxy$ như hình vẽ.

\begin{align*}
   \vec{F} &= \vec{F}_1 + \vec{F}_2 + \vec{F}_3 \\
   &= F_1 \vec{j} + F_2 \cos(30^\circ) \vec{i} + F_2 \sin(30^\circ) \vec{j} + F_3 \cos(210^\circ) \vec{i} + F_3 \sin(210^\circ) \vec{j} \\
   &= 20 \vec{j} + 30 \cdot \frac{\sqrt{3}}{2} \vec{i} + 30 \cdot \frac{1}{2} \vec{j} + 40 \cdot \left(-\frac{\sqrt{3}}{2}\right) \vec{i} + 40 \cdot \left(-\frac{1}{2}\right) \vec{j} \\
   &= -5\sqrt{3} \vec{i} + 15 \vec{j}.
\end{align*}

Vậy, độ lớn hợp lực là:
\begin{equation*}
   \left|\vec{F}\right| = \sqrt{(-5\sqrt{3})^2 + 15^2} = 10\sqrt{3} \approx \boxed{\qty{17.3}{N}}.
\end{equation*}